% SOURCE_AUTHORITY: sga5_fr_workpass.tex, lines 12943--12969 (LF-normalized unit).
% SOURCE_SHA256: 791F4EFFC5E02832D5D77ED03518C8156D6F07E4C8238B03545DB93D883FBB28
a) Seja $\Psi$ um complexo limitado à esquerda de feixes de $\Lambda[G]$-módulos sobre $C$. Tem-se um isomorfismo
\[
\R\Gamma^G(\RG_C(\Psi))\simeq\RG_C(\R\underline{\Gamma}^G(\Psi)),
\tag{7.3}
\]
onde $\underline{\Gamma}^G$ é o funtor «subfeixe das seções $G$-invariantes». De fato, $\Gamma_C$ transforma um feixe injetivo de $\Lambda[G]$-módulos em um $\Lambda[G]$-módulo injetivo e $\underline{\Gamma}^G$ transforma um feixe injetivo de $\Lambda[G]$-módulos em um feixe injetivo de $\Lambda$-módulos (utilizem-se, para isso, por exemplo, adjuntos evidentes de $\Gamma_C$ e $\underline{\Gamma}^G$); por outro lado, tem-se
\[
\Gamma^G\circ\Gamma_C=\Gamma_C\circ\underline{\Gamma}^G.
\]
Tomando os funtores derivados de ambos os membros e aplicando ([5] prop. 3.1), obtém-se, portanto, (7.3).

b) Seja $\Phi$ um complexo limitado à esquerda de feixes de $\Lambda$-módulos sobre $C$, nulo fora de $U$. Tem-se um isomorfismo
\[
\Phi\simeq \R\underline{\Gamma}^G(p_*(\Phi')),
\qquad \Phi'=p^*(\Phi).
\tag{7.4}
\]
Observe-se primeiro que os funtores $p^*$, $\underline{\Gamma}^G\circ p_*$ estabelecem equivalências quase inversas entre a categoria dos feixes de $\Lambda$-módulos sobre $C'$ sobre os quais opera $G$ e que são nulos fora de $U'=p^{-1}(U)$, e a categoria dos feixes de $\Lambda$-módulos sobre $C$ nulos fora de $U$. Tem-se, portanto, o isomorfismo
\[
\Phi\simeq \R\bigl((\underline{\Gamma}^G\circ p_*)\circ p^*\bigr)(\Phi).
\]
Por outro lado, pode-se aplicar novamente ([5] Prop. 3.1), que fornece, tendo em conta que $\R p^*=p^*$, $\R p_*=p_*$:
\[
\Phi\simeq \R(\underline{\Gamma}^G\circ p_*)(p^*(\Phi))
\simeq (\R\underline{\Gamma}^G)(p_*p^*(\Phi)),
\]
que é o isomorfismo procurado.
